\SourceLeaf{377}%
nale bedrohen, als die Beobachtungen selbst in hohem Grade erschweren. Man fand sich durch diese Umstände bewogen, eine verhältnissmässig grosse Anzahl ziemlich kleiner Dreiecke zu bilden: es sind sechszehn, und die kleinste Verbindungsseite ist nur 18671 Meter lang. Wir dürfen jedoch nicht unbemerkt lassen, dass die Heliotrope, welche alle Signale ganz entbehrlich, und die Messung der Winkel in den allergrössten Dreiecken eben so leicht und scharf, wie bei dem kleinsten, machen, damals in Italien noch nicht bekannt waren.

Zur Messung der Winkel dienten achtzollige Theodolithen von \textsc{Reichenbach}. Die Piemontesischen und Oesterreichischen Officiere theilten sich nicht in die Arbeit, sondern jene und diese bestimmten sämmtliche Winkel des Systems unabhängig für sich. Man erhielt also von jedem einzelnen Winkel zwei Bestimmungen, aus denen nach Massgabe der Anzahl der Serien, die dazu concurrirt hatten, das Mittel als Definitivwerth angenommen wurde. Meistens beruhen die Resultate der Piemontesischen Officiere auf sechs Serien, jede zu 10 Repetitionen; die der Oesterreichischen grösstentheils auf zwei, einige auf drei oder vier Serien. Alle Messungen sind im grössten Detail abgedruckt, doch ohne Nennung der Beobachter, von denen jede einzeln herrührt.

Bei einer so ausgedehnten Operation hat die Kenntniss der bei den Winkelbestimmungen erreichten Genauigkeit ein grosses Interesse. Die Winkelsummen in den einzelnen Dreiecken bieten ein Mittel dazu dar, welches freilich nach Umständen etwas trüglich sein kann. Darf man die vorliegenden danach beurtheilen, so haben sie allerdings eine bewunderungswürdige Genauigkeit. Der grösste Fehler der Winkelsumme bei den 16 Dreiecken ist nur \(1''16\); der mittlere Fehler findet sich \(0''70\), und der mittlere Fehler einzelner Winkel würde folglich nur \(0''40\) sein. Prüfungsmittel durch Diagonalrichtungen oder Polygonbildungen sind gar nicht vorhanden. Allein die Vergleichung der doppelten Bestimmungen der 48 Winkel unter sich deutet ganz entschieden auf eine bei weitem grössere Ungenauigkeit der Resultate hin; wir finden hier 13 wo der Unterschied über \(3''\), und darunter 5 wo er über \(5''\) steigt, ja bei einer, gleich in dem ersten Dreiecke, weicht die auf 80 Repetitionen gegründete Bestimmung der Piemontesischen Officiere von der auf 48 Repetitionen beruhenden der Oesterreichischen um \(9''2\) ab. Bei so grossen Differenzen kann man sich der Vermuthung nicht erwehren, dass die richtige Würdigung der eigentlichen Genauigkeit der Messungen noch von Nebenumständen abhängt, von welchen das Werk uns keine Kenntniss gibt.
